% META: {"id": "TEXP-ARI-EX-015", "chapitre": "TEXP-ARITHMETIQUE", "type_objet": "exercice", "capacites": ["TEXP-ARI-C3"], "capacites_codes": ["C3"], "methodes": ["M3"], "parcours": 2, "competences": ["calculer", "raisonner"], "duree_min": 15, "mode_creation": "ex_nihilo", "sources_inspiration": [], "fichier_tex": "chapitres/TEXP-ARITHMETIQUE/exercices/TEXP-ARI-EX-015.tex", "corrige_tex": "chapitres/TEXP-ARITHMETIQUE/corriges/TEXP-ARI-CO-015.tex", "status": "generated", "gestes": ["calcul", "traiter_cas_limite"]}
% BEGIN-VERIFY
% import math
% def est_premier(n):
%     if n < 2:
%         return False
%     for p in range(2, int(math.isqrt(n)) + 1):
%         if n % p == 0:
%             return False
%     return True
% assert math.isqrt(2027) == 45
% assert est_premier(2027) is True
% assert 2021 == 43 * 47 and est_premier(2021) is False
% assert 2027 % 3 == 2 and 2027 % 7 == 4 and 2027 % 43 == 6
% END-VERIFY
\begin{exercice}{TEXP-ARI-EX-015}{2}{15}
\begin{enumerate}
  \item Justifier qu'il suffit de tester les diviseurs premiers $p\leq\sqrt{n}$
        pour décider si $n$ est premier.
  \item Le nombre $2\,027$ est-il premier ? Détailler les tests effectués.
  \item Même question pour $2\,021$.
\end{enumerate}
\end{exercice}
